\section{approach}
\label{sec:approach}

\subsection{Overview Design of \imss}
\begin{figure}[H]
    \center
    \includegraphics[scale=0.45]{imss}
    \caption{\label{fig:imss}Schematic view of \imss.}
\end{figure}
\imss finds a string that satisfies constraints specified in the input,
or decides that no satisfying string exists. \imss utilizes three different
string solvers: \hampi, \dprle, \zstr to solve string
constraints. In summary, \imss works in the following four steps (Figure~\ref{fig:imss}):
\begin{enumerate}
    \item Parse the input constraints into an abstract syntax tree (AST)
    representation of the \imss language.
    \item Translate the AST into the formats that corresponds to the three different solvers.
    \item Execute all the three solvers in parallel.
    \item Gather all the results. If none of the solvers can solve the constraints
    specified in the input, report that the constraints are incompatitble.
\end{enumerate}
The input format to \imss is specified by the grammar in Grammar~\ref{fig:grammar}.

\begin{Grammar}[H]
    \setlength{\grammarindent}{5em} % increase separation between LHS/RHS
    \begin{grammar}
        <Input> ::= <Statement>*                    \hfill \imss input

        <Statement> ::= String(<ID>)                \hfill string
                \alt Int(<ID>)                      \hfill int
                \alt StringDecl(<ID>, <STR_LIT>)    \hfill string witn value
                \alt Regex(<ID>, <RegStm>)          \hfill regex lang.
                \alt CFG(<ID>, <CFGStm>)            \hfill CFG lang.
                \alt <AssertStm>                    \hfill assert statement

        <RegStm> ::= <STR_LIT>                      \hfill string literal
                \alt <ID>                           \hfill variable name
                \alt Fix(<ID>, <INT_LIT>)           \hfill CFG fixed-sizing
                \alt Or(<RegStm>*)                  \hfill union
                \alt Concat(<RegStm>*)              \hfill concatenation
                \alt Star(<RegStm>)                 \hfill Kleene star

        <AssertStm> ::= AssertIn(<ID>,<ID>)         \hfill membership
                \alt Equal(<Exp>,<Exp>)             \hfill string equality
                \alt Contains(<StrExp>,<StrExp>)    \hfill string contains
                \alt LessThan(<IntExp>,<IntExp>)    \hfill integer comparison

        <Exp> ::= <IntExp>                          \hfill integer expression
                \alt <StrExp>                       \hfill string expression
                \alt <RegExp>                       \hfill regex expression
                \alt <ID>                           \hfill variable name

    \end{grammar}
 \caption{\label{fig:grammar}Subset of \imss's input language, shown in BNF grammars.
 Non-terminals are shown in angle brackets. \imss's input consists of a list of statements.
 Statements can be variable declaration or assertion statements which represent
 string constraints. Some non-terminals are omitted for simplicity.}
 \end{Grammar}

\subsection{Example}

In order to demonstrate the four steps of \imss, consider the example below:
\renewcommand\lstlistingname{Code}
\begin{lstlisting}[label={code:example},frame=bt,numbers=none,basicstyle=\ttfamily\small,
language=C,captionpos=b,caption={Input constraint in \imss format}]
1.  String(v);
2.  RangedLength(v, 1, 100);
3.  Contains(v, "abc");
4.  Regex(r, Star(Or("ab", "cd")))
5.  AssertIn(v, r);
6.  SolveAll();
\end{lstlisting}
This is a constraint for a string that contains the substring $abc$, and
the string is a member of the regular expression, $(ab|cd)*$.
One of the strings which satisfy this constraint is ``abcd''. The string can be of any length
between 1 and 100. \imss will parse the input into an AST (Figure~\ref{fig:exp-ast}) according to Grammar~\ref{fig:grammar}.

\begin{figure*}
    \centering
    \includegraphics[scale=0.65]{example-ast}
    \caption{\label{fig:exp-ast}Abstract syntax tree of the example Code~\ref{code:example}.}
\end{figure*}

Then, \imss will translate the AST into all three input formats by traversing the tree and
translating each tree node into the string format for each language. If a language does not support
some functionality (\imss cannot translate some node into that specific language),
then translation to that language will fail and \imss will output ``incompatible'' for that language. After the translations, \imss executes all of the generated input formats
on its corresponding solver in parallel and outputs the fastest solution (or all solutions depending on the configuration for \imss).
